% ============================================================
%  style.tex —— 通用期末模拟卷版式（试卷册）
%  依赖：meta.tex（课程元信息）、ctexart 文档类、xelatex 编译
%  设计目标：无课程耦合，任何课程的试卷都复用这一份
% ============================================================

% ============================================================
%  meta.tex —— 课程元信息（整份试卷唯一需要改动的配置）
%  换一门课，只改这个文件，style.tex 不用动
% ============================================================

% ---- 学校与课程 ----
\newcommand{\SchoolName}{××大学}
\newcommand{\CollegeName}{}          % 学院名，留空则不打印；填了会接在学校名后面
\newcommand{\CourseName}{信号与系统}
\newcommand{\TermName}{2026—2027 学年\quad 第 1 学期}
\newcommand{\ExamDate}{}             % 考试日期，留空则不打印

% ---- 卷面抬头 ----
\newcommand{\MajorList}{电子信息工程、通信工程、自动化}   % 适用专业
\newcommand{\ExamForm}{闭卷}                              % 开卷 / 闭卷 / 半开卷
\newcommand{\ExamTime}{120 分钟}
\newcommand{\TotalPoints}{100}
\newcommand{\ScopeNote}{全课程}                           % 命题范围

% ---- 卷面要素开关 ----
% 1 = 显示，0 = 隐藏。密封线在双面印刷或不要求装订时可关掉
\newcommand{\ShowSealLine}{1}
\newcommand{\ShowScoreTable}{1}

% ---- 密封线文字（仅 \ShowSealLine=1 时生效）----
\newcommand{\SealText}{%
班级\underline{\hspace{1.5cm}}\quad 学号\underline{\hspace{1.75cm}}\quad 姓名\underline{\hspace{1.45cm}}\qquad
{\heiti 密}\ {\heiti 封}\ {\heiti 线}\qquad 请勿在密封线内答题}

% ---- 试卷说明（三到四条为宜，超出会挤掉正文空间）----
\newcommand{\ExamNotice}{%
1. 答案一律写在试卷相应位置，写在草稿纸上无效；%
2. 计算题须写出主要公式与演算过程，只有结果不给分；%
3. 答卷前请将密封线内项目填写清楚。}

% ---- 答案册评分说明 ----
\newcommand{\GradeNotice}{%
填空题按空给分，答案正确即得分；选择题、判断题不倒扣分；%
简答题按要点给分，答出关键术语与结论即可，表述完整、条理清楚者可在该题分值内适当调整；%
计算题按步骤给分，公式正确而结果错误的扣 1～2 分，只有结果无过程的不得分。}


% ---------- 版心：左侧留出 1.6cm 密封线，故左边距偏大 ----------
\usepackage[a4paper,left=2.7cm,right=1.9cm,top=2.1cm,bottom=1.9cm,
            headsep=0.35cm,footskip=0.9cm]{geometry}
\usepackage{amsmath,amssymb}
\usepackage{graphicx}
\usepackage{array,tabularx,multirow}
\usepackage{enumitem}
\usepackage{xcolor}
\usepackage{tikz}
\usetikzlibrary{arrows.meta,positioning,calc,patterns}
\usepackage{eso-pic}
\usepackage{fancyhdr}
\usepackage{lastpage}
\usepackage{needspace}
\usepackage{listings}
\usepackage{etoolbox}   % \ifdefempty：让卷头可选行留空就不打印

% ---------- 中文字体：有则用宋/黑/楷，无则回退 ctex 默认（macOS/Linux 友好）----------
\IfFontExistsTF{SimSun}{%
  \setCJKmainfont{SimSun}[BoldFont=SimHei,ItalicFont=KaiTi]%
  \setCJKsansfont{SimHei}%
}{}
\IfFontExistsTF{FangSong}{\setCJKmonofont{FangSong}}{}
\IfFontExistsTF{Consolas}{\setmonofont{Consolas}}{}

% ---------- 全局行距与段距：行距偏紧以容纳更多题目 ----------
\linespread{1.32}
\setlength{\parindent}{0pt}
\setlength{\parskip}{0pt}
% 中文断行遇到"、"连排时容易挤出边界，给一点应急伸缩量吸收 2–6pt 的溢出
\setlength{\emergencystretch}{2em}

% ---------- 代码块（编程类课程用；不写代码的课程留空即可）----------
\IfFontExistsTF{Consolas}{\setmonofont{Consolas}}{}
\lstdefinestyle{excode}{
  basicstyle=\ttfamily\fontsize{7.6pt}{10.4pt}\selectfont,
  columns=fullflexible,
  keepspaces=true,
  frame=none,
  aboveskip=1pt,belowskip=1pt,
  xleftmargin=1.4em,
  showstringspaces=false,
  breaklines=false,
  extendedchars=false
}
\lstnewenvironment{code}[1][]{\lstset{style=excode,#1}}{}

% ============================================================
%  页眉页脚
% ============================================================
\newcommand{\juanmark}{A}
\pagestyle{fancy}
\fancyhf{}
\renewcommand{\headrulewidth}{0.4pt}
\fancyhead[L]{\small\songti \SchoolName《\CourseName》期末模拟试卷（\juanmark\ 卷）}
\fancyhead[R]{\small\songti 第 \thepage\ 页\quad 共 \pageref{LastPage} 页}
\fancyfoot[C]{}

% ============================================================
%  卷头
%  \juantou{A}  —— 参数为卷别标记
%  学校/学院/课程/学期/专业/日期全部从 meta.tex 取；留空的行自动不打印
% ============================================================
\newcommand{\juantou}[1]{%
\clearpage   % 关键：先输出上一卷最后一页，此时 \juanmark 还是旧值，页眉才不会串卷
\renewcommand{\juanmark}{#1}%
\begin{center}
{\zihao{5}\songti \SchoolName\ifdefempty{\CollegeName}{}{\ \CollegeName}\ \ \TermName}\\[4pt]
{\zihao{3}\heiti 《\CourseName》期末考试模拟试卷（#1 卷）}\\[2pt]
\ifdefempty{\MajorList}{}{\zihao{6}\kaishu 适用专业：\MajorList\\[3pt]}%
{\zihao{6}\songti 考试形式：\ExamForm\ \ \ 考试时间：\ExamTime\ \ \ 满分：\TotalPoints\ 分\ \ \ 命题范围：\ScopeNote}%
\ifdefempty{\ExamDate}{}{\\[2pt]{\zihao{6}\songti 考试日期：\ExamDate}}%
\end{center}
\vspace{-2pt}
\hrule height 1.2pt
\vspace{3pt}
{\zihao{6}\songti 试卷说明：\ExamNotice}
\vspace{4pt}
}

% ============================================================
%  得分表
%  大题数固定 5 或 6，用对应版本；分值之和自动计算
%  \defenBiao{15}{20}{10}{20}{25}{10}      % 六个大题
%  \defenBiaoFive{15}{20}{20}{25}{20}      % 五个大题
% ============================================================
\ExplSyntaxOn
% 把 "15,20,10" 求和，用于得分表总分自动计算，杜绝手算错分
\NewDocumentCommand{\ScoreSum}{m}{\int_eval:n {\clist_use:nn {#1} {+}}}
\ExplSyntaxOff

\newcommand{\defenBiao}[6]{%
\ifnum\ShowScoreTable=1
\begin{center}
\zihao{6}
\begin{tabularx}{\textwidth}{|>{\centering\arraybackslash}p{1.9cm}|*{6}{>{\centering\arraybackslash}X}|>{\centering\arraybackslash}p{1.5cm}|>{\centering\arraybackslash}p{1.6cm}|>{\centering\arraybackslash}p{1.6cm}|}
\hline
题\ \ 号 & 一 & 二 & 三 & 四 & 五 & 六 & 总分 & 合分人 & 复核人 \\ \hline
满\ \ 分 & #1 & #2 & #3 & #4 & #5 & #6 & \ScoreSum{#1,#2,#3,#4,#5,#6} & & \\ \hline
得\ \ 分 & & & & & & & & & \\ \hline
\end{tabularx}
\end{center}
\vspace{2pt}
\fi}

\newcommand{\defenBiaoFour}[4]{%
\ifnum\ShowScoreTable=1
\begin{center}
\zihao{6}
\begin{tabularx}{\textwidth}{|>{\centering\arraybackslash}p{1.9cm}|*{4}{>{\centering\arraybackslash}X}|>{\centering\arraybackslash}p{1.5cm}|>{\centering\arraybackslash}p{1.6cm}|>{\centering\arraybackslash}p{1.6cm}|}
\hline
题\ \ 号 & 一 & 二 & 三 & 四 & 总分 & 合分人 & 复核人 \\ \hline
满\ \ 分 & #1 & #2 & #3 & #4 & \ScoreSum{#1,#2,#3,#4} & & \\ \hline
得\ \ 分 & & & & & & & \\ \hline
\end{tabularx}
\end{center}
\vspace{2pt}
\fi}

\newcommand{\defenBiaoFive}[5]{%
\ifnum\ShowScoreTable=1
\begin{center}
\zihao{6}
\begin{tabularx}{\textwidth}{|>{\centering\arraybackslash}p{1.9cm}|*{5}{>{\centering\arraybackslash}X}|>{\centering\arraybackslash}p{1.5cm}|>{\centering\arraybackslash}p{1.6cm}|>{\centering\arraybackslash}p{1.6cm}|}
\hline
题\ \ 号 & 一 & 二 & 三 & 四 & 五 & 总分 & 合分人 & 复核人 \\ \hline
满\ \ 分 & #1 & #2 & #3 & #4 & #5 & \ScoreSum{#1,#2,#3,#4,#5} & & \\ \hline
得\ \ 分 & & & & & & & & \\ \hline
\end{tabularx}
\end{center}
\vspace{2pt}
\fi}

% ============================================================
%  大题标题（自带得分/评阅人格）
%  \dati{一、填空题}{每空 1 分，共 15 分}
% ============================================================
\newcommand{\dati}[2]{%
\needspace{4.6cm}
\vspace{8pt}
\begin{center}\zihao{6}
\begin{tabularx}{\textwidth}{|>{\raggedright\arraybackslash}p{0.635\textwidth}|>{\centering\arraybackslash}X|>{\centering\arraybackslash}X|}
\hline
\heiti #1（#2） & 得\ 分 & 评\ 阅\ 人 \\ \hline
 & \rule{0pt}{1.0cm} & \\ \hline
\end{tabularx}
\end{center}
\vspace{1pt}}

% ============================================================
%  题目与答题区
% ============================================================
% 单行题目：自动悬挂缩进，续行与小序号对齐
\newcommand{\xiaoti}[1]{\par\vspace{3.2pt}\noindent\hangindent=1.6em\hangafter=1 #1\par}

% 填空线：\kong 默认 2.4cm，可按答案长度传参 \kong[1.3cm]
\newcommand{\kong}[1][2.4cm]{\underline{\hspace{#1}}}

% 答题空白：用不可分割的 parbox，避免被排版吞掉或跨页撕裂
\newcommand{\kongbai}[1]{\par\vspace{0.12cm}\noindent\parbox[t][#1][t]{\linewidth}{\rule{0pt}{0pt}}\par}

% 题目+答题区整体不跨页：\begin{ti} ... \kongbai{4cm} ... \end{ti}
\newenvironment{ti}{\par\vspace{3pt}\noindent\begin{minipage}{\linewidth}%
\setlength{\parindent}{0pt}\setlength{\parskip}{0pt}}{\end{minipage}\par}

% 插图题注（统一小一号字，居中）
\newcommand{\tucaption}[1]{\begin{center}\zihao{6}\songti #1\end{center}}

% ============================================================
%  密封线（每页背景层，可用 \ShowSealLine 关闭）
% ============================================================
\newcommand{\sealLine}{%
\ifnum\ShowSealLine=1
\AddToShipoutPictureBG{%
\begin{tikzpicture}[remember picture,overlay]
\draw[densely dashed,line width=0.5pt]
  ([xshift=1.62cm]current page.north west)--([xshift=1.62cm]current page.south west);
\node[rotate=-90,font=\zihao{6}\songti]
  at ([xshift=1.02cm,yshift=-14.85cm]current page.north west) {\SealText};
\end{tikzpicture}}%
\fi}
